% Standalone section -- paste into the weekly report, or \input it.
% Week 3 working note: what a centerline is, how a 1D curve bends and twists, the established
% tortuosity measures (including the curvature-squared integral), and how representation learning
% could find tortuosity structure in the tree without being told what tortuosity is.
%
% PROVENANCE OF THE NUMBERS. Every case-1 number below was printed by the repo's own scripts on
% images/1 (ImageCAS-X case 1) on 2026-09-11:
%   uv run python scripts/discover_centerline.py images/1           (graph, spacing, runs)
%   uv run python scripts/discover_tortuosity.py images/1 --sweep   (metrics, sigma sweep)
% One case, so they illustrate; they are not population statistics.
%
% READING DISCIPLINE.
%   - kashyap2022tortuosity: the four metric definitions and R^2 values were read from the full
%     text (PMC8764056, Tables 1 and 4).
%   - korona2025tracing: read in the PDF in Former_students_work/ (sections 2 and 5.2).
%   - Every other new reference (hart1999tortuosity ... higgins2017betavae) was read at abstract
%     level only. What is attributed to them is what their abstracts or search-result summaries
%     say, and the grisan2008tortuosity formula is from memory. Re-read before the thesis.
%   - The noise-bias estimate in the T_kappa subsection is my own back-of-envelope derivation,
%     not a literature result. The rough estimate first given in chat (6 eps^2/h^4) undercounted
%     by a factor of 2; the text uses the corrected 12 eps^2/h^4, which is of the same order.
%
% Citation keys resolve against thesis/refs.bib. The block "added 2026-09-11" was created for
% this file, so refs_clean.bib has to be regenerated after editing:
%   python scripts/make_refs_clean.py
%
% CITATION WANTED: the fundamental theorem of space curves (any differential-geometry textbook,
% e.g. do Carmo) and the Dean number (Dean 1927) are stated without citations.

\section*{Tortuosity from centerlines: measures and learned representations}

This note gathers the week's thinking on tortuosity in one place. It starts from what a centerline
actually contains, because every measure below is computed from it. Then it covers how a curve
that is one-dimensional can still bend and twist, and what the tortuosity of a real coronary
centerline is made of. It then turns to the established hand-crafted measures and, last, to
whether a neural network could find the same structure --- or more --- on its own.

% ---------------------------------------------------------------------------------------------
\subsection*{1. What a centerline is}

A centerline is the curve that runs through the middle of a vessel: the tube reduced to its axis.
The standard definition, the one VMTK implements, is the path traced by the centres of the largest
spheres that fit inside the lumen, the radius of each sphere being the local vessel radius
\cite{antiga2008vmtk}. ``The middle'' is well defined only where the vessel is a clean tube; at
bifurcations and stenoses different methods put it in different places.

\textbf{How it is made.} Three routes are common. A segmentation mask can be thinned to a
one-voxel skeleton and then smoothed: this is how the ImageCAS-X centerlines were produced, from
the manually corrected masks with Gaussian smoothing at $\sigma = 0.5$~mm
\cite{bransby2026imagecasx}. A surface mesh can be reduced to its medial axis (VMTK). Or a tracker
can walk the vessel in the image directly. The Korona--Baldachowski tracer is of this last kind:
a CNN scores the 26 neighbouring voxels on a 0.5~mm grid and the tracer steps to the best one
\cite{korona2025tracing}.

\textbf{What the file holds.} A centerline in ImageCAS-X is VTK PolyData, one file per side. For
case 1, the left tree (\Cref{tab:week3-centerline-file}) holds:
\begin{itemize}
  \item \textbf{Points}: an $857 \times 3$ float32 array of $x, y, z$ in mm, in the scanner frame
    shared with the CT and the mask. Median spacing is 0.474~mm (range 0.342--0.986~mm), about one
    point per voxel.
  \item \textbf{Cells}: 13 polylines, each a list of point indices. These are the edges of the
    tree, from one branch or end point to the next. The order of the point array itself carries no
    meaning, and the cells carry no flow direction; a proximal-to-distal ordering has to be
    recovered by a traversal rooted at the ostium.
  \item \textbf{Per-point attributes}: an integer \texttt{segment\_label} and a string
    \texttt{segment\_name} (LM, LAD, LCx, D1, \ldots), and three 0/1 flags for the ostium
    (\texttt{start\_points}, 1 point), the bifurcations (\texttt{branch\_points}, 6) and the
    distal tips (\texttt{end\_points}, 8).
\end{itemize}
The right tree has the same layout, with 500 points and 5 edges.

\begin{table}[htbp]
  \centering
  \small
  \begin{tabular}{@{}llp{0.46\linewidth}@{}}
    \toprule
    Array & Type & Meaning \\
    \midrule
    points & float32, $N \times 3$ & $x, y, z$ in mm, scanner frame \\
    cells (polylines) & lists of point indices & edges of the tree, branch point to branch point \\
    \texttt{segment\_label} & int & anatomical segment: 1 = LM, 2 = LAD, 3 = LCx, \ldots,
      14 = Other \\
    \texttt{segment\_name} & string & the same, as text \\
    \texttt{start\_points} & 0/1 & ostium \\
    \texttt{branch\_points} & 0/1 & bifurcations \\
    \texttt{end\_points} & 0/1 & distal tips \\
    \bottomrule
  \end{tabular}
  \caption{Everything stored in an ImageCAS-X centerline file. No radius, tangent, curvature or
    tortuosity is stored: all of it is derived from the coordinates and the connectivity.}
  \label{tab:week3-centerline-file}
\end{table}

\textbf{One LAD in numbers.} In case 1 the LAD has 309 points, so $309 \times 3 = 927$ numbers
describe the whole vessel. It is stored in four pieces, cut wherever a side branch leaves it:
18 points (7.2~mm) from the LM split to the D1 origin, 51 points (22.8~mm) to D2, 34 points
(15.9~mm) to the next branch, and 210 points (106.3~mm) to the tip --- about 152~mm in all.
Consecutive points lie a fraction of a millimetre apart; the shape of the vessel is in how those
small steps change direction.

This has a direct consequence for the existing script. Its per-\emph{segment} rows measure only the
longest contiguous piece of each label, so the ``LAD'' row below is the distal 106~mm, not the whole
vessel. A whole-LAD tortuosity needs the four pieces chained from the LM split to the tip. The
ostium-to-tip paths already do that, but they include the left main as well.

\textbf{What a centerline is not.} It holds no width and no lumen shape. Those come from the mask
or the surface mesh, or from a radius attached to each point; the script's radius table gets it
from the distance between each centerline point and the surface.

% ---------------------------------------------------------------------------------------------
\subsection*{2. How a one-dimensional curve bends and twists}

``One-dimensional'' describes the curve, not the space it sits in. A single number --- the arc
length $s$ from the ostium --- says where along it you are, but each point has three coordinates,
so the curve can bend in any direction, the way a wire can be bent into a spring. With $\mathbf r(s)$
the position and $\mathbf T = \mathrm d\mathbf r/\mathrm ds$ the unit tangent, two quantities
describe everything it does:
\begin{itemize}
  \item \textbf{Curvature} $\kappa(s) = \lvert \mathrm d\mathbf T/\mathrm ds \rvert$, in
    mm$^{-1}$: how fast the direction turns, i.e.\ how sharply it bends. A bend of radius $R$ has
    $\kappa = 1/R$.
  \item \textbf{Torsion} $\tau(s)$: how fast the \emph{plane} of the bend rotates about the
    tangent, i.e.\ how much the curve leaves the plane it was bending in. This is the twist.
\end{itemize}
For a curve sampled at arbitrary speed, which is what finite differences give,
\begin{equation}
  \kappa = \frac{\lvert \mathbf r' \times \mathbf r'' \rvert}{\lvert \mathbf r' \rvert^{3}},
  \qquad
  \tau = \frac{(\mathbf r' \times \mathbf r'') \cdot \mathbf r'''}{\lvert \mathbf r' \times
  \mathbf r'' \rvert^{2}},
  \label{eq:week3-frenet}
\end{equation}
which is what \texttt{scripts/discover\_tortuosity.py} computes.

\textbf{The discrete picture.} Three consecutive points define a plane, and the angle between
their two steps is the bend. A fourth point either stays in that plane ($\tau = 0$) or lifts out of
it; the dihedral angle between the planes through points (1,2,3) and (2,3,4) is the twist.
\Cref{tab:week3-curves} lists the standard cases.

\begin{table}[htbp]
  \centering
  \small
  \begin{tabular}{@{}lll@{}}
    \toprule
    Curve & $\kappa$ & $\tau$ \\
    \midrule
    straight line & 0 & undefined \\
    circle, or any bend in one plane & $> 0$ & 0 \\
    helix (a spring) & constant $> 0$ & constant $\neq 0$ \\
    S-bend in one plane & $> 0$, side flips at inflection & 0 \\
    S-bend whose second bend leaves the plane & $> 0$ & $\neq 0$ where it leaves \\
    \bottomrule
  \end{tabular}
  \caption{Curvature and torsion of standard curves.}
  \label{tab:week3-curves}
\end{table}

\textbf{Nothing is lost by treating the curve as 1D.} By the fundamental theorem of space curves,
$\kappa(s)$ and $\tau(s)$ along the whole length determine the curve uniquely up to its position
and orientation. Its bending and twisting are entirely contained in those two profiles.

A coarser measure of twist is \emph{non-planarity}: the fraction of the curve's spread that lies
outside its best-fitting plane (the smallest eigenvalue of the point covariance over their sum).
In case 1 the LM is essentially flat (0.0002), the LAD's distal piece reaches 0.009, and the RCA is
the most three-dimensional at 0.022, as it wraps around the heart.

Two caveats. Torsion is a third derivative, noisier still than curvature, and on a nearly straight
stretch the bending plane is undefined, so $\tau$ is usually trusted only where $\kappa$ is clearly
above zero, and only after smoothing. And ``twist'' can mean two things: the centerline gives the
\emph{path} twisting out of plane, which is what drives helical flow; it cannot tell whether the
\emph{tube} rotates about its own axis, which for a round lumen has no visible effect.

% ---------------------------------------------------------------------------------------------
\subsection*{3. What a centerline's tortuosity is made of}

A centerline certainly has tortuosity; every measure below is defined on it. Case 1 shows plenty
(\Cref{tab:week3-case1}). The question is which part of it is wanted, because the number mixes
three things.

\begin{table}[htbp]
  \centering
  \small
  \begin{tabular}{@{}lrrrrrr@{}}
    \toprule
    Segment & Arc (mm) & Chord (mm) & $L/D$ & SOAM (rad/mm) & $\kappa_{\max}$ (mm$^{-1}$) &
    Non-planarity \\
    \midrule
    LM  &   5.9 &  5.8 & 1.009 & 0.088 & 0.121 & 0.0002 \\
    LAD & 104.5 & 79.6 & 1.313 & 0.072 & 0.346 & 0.0087 \\
    LCx &  43.3 & 36.1 & 1.202 & 0.086 & 0.234 & 0.0022 \\
    RCA & 116.0 & 47.9 & 2.424 & 0.100 & 0.413 & 0.0216 \\
    \bottomrule
  \end{tabular}
  \caption{Case 1 at $\sigma = 1$~mm. Each row is the longest contiguous piece of that label
    (the LAD row is its distal 106~mm). SOAM is the sum of turning angles per mm, which for
    finely sampled curves equals the mean absolute curvature (0.072~mm$^{-1}$ for the LAD).}
  \label{tab:week3-case1}
\end{table}

\begin{enumerate}
  \item \textbf{Extraction noise.} Skeletonising voxels leaves jitter at the voxel scale.
    \Cref{tab:week3-sweep} shows it: smoothing the LAD from $\sigma = 0$ to 1~mm drops the mean
    curvature from 0.172 to 0.072~mm$^{-1}$ and the inflection count from 43 to 4, while $L/D$
    barely moves (1.328 to 1.313). Real bends survive smoothing: at $\sigma = 2$~mm the LAD still
    has $\kappa_{\max} = 0.146$~mm$^{-1}$, a bend radius of about 7~mm.
  \item \textbf{The heart's own shape.} The RCA's $L/D = 2.4$ comes mostly from the vessel wrapping
    around the atrioventricular groove, following the heart's surface. Clinically ``tortuous''
    means extra bends on top of that course, and a whole-vessel $L/D$ cannot separate the two.
    Local measures mostly can, because the heart's own bend is gentle (a radius of about 40~mm is
    $\kappa \approx 0.025$~mm$^{-1}$).
  \item \textbf{The vessel's own local bends}, sharper than the heart's curve and broader than the
    voxel noise. This is the part that matters for flow and for where plaque starts.
\end{enumerate}

\begin{table}[htbp]
  \centering
  \small
  \begin{tabular}{@{}lrrrrrrrr@{}}
    \toprule
    & \multicolumn{4}{c}{LAD (distal piece)} & \multicolumn{4}{c}{RCA} \\
    \cmidrule(lr){2-5}\cmidrule(l){6-9}
    $\sigma$ (mm) & $L/D$ & infl. & $\bar\kappa$ & $\kappa_{\max}$ & $L/D$ & infl. & $\bar\kappa$ &
    $\kappa_{\max}$ \\
    \midrule
    0.0 & 1.328 & 43 & 0.172 & 1.133 & 2.475 & 43 & 0.217 & 1.048 \\
    0.5 & 1.319 & 14 & 0.116 & 0.632 & 2.451 & 24 & 0.152 & 0.664 \\
    1.0 & 1.313 &  4 & 0.072 & 0.346 & 2.424 &  7 & 0.101 & 0.413 \\
    2.0 & 1.302 &  0 & 0.052 & 0.146 & 2.368 &  1 & 0.080 & 0.200 \\
    \bottomrule
  \end{tabular}
  \caption{Case 1: how the metrics move with the smoothing scale $\sigma$ applied before
    differentiation (curvature in mm$^{-1}$). $L/D$ is a ratio of lengths and barely moves;
    everything derivative-based does.}
  \label{tab:week3-sweep}
\end{table}

What a centerline cannot say at all: the radius (the same bend disturbs flow more in a wider
vessel, which is what the Dean number expresses), the curvature of the inner versus the outer wall,
and the heartbeat --- coronaries move and bend with each beat, so a CT centerline is a snapshot at
the reconstructed phase.

% ---------------------------------------------------------------------------------------------
\subsection*{4. Established tortuosity measures}

The measures are grouped by the derivative of the centerline they need, because on discretised
centerlines each extra derivative amplifies noise.

\textbf{A. Length ratios (positions only).}
The \emph{tortuosity index} or distance metric, $L/D$, is arc length over chord. It is very stable,
but blind to shape: all six curves in \Cref{fig:week3-tortuosity-ambiguity} share $L/D = 1.25$.
It was judged of little use in 3D intracerebral vessels \cite{bullitt2003tortuosity}, and showed no
relation to low wall shear stress in 127 CAD-free left main bifurcations \cite{kashyap2022tortuosity}.
The \emph{tortuosity density} of Grisan et al.\ \cite{grisan2008tortuosity} cuts the curve at its
$n-1$ inflections into $n$ pieces of constant bending side, with arc $L_i$ and chord $C_i$, and
combines
\begin{equation}
  \tau_d = \frac{n-1}{n}\,\frac{1}{L}\sum_{i=1}^{n}\Bigl(\frac{L_i}{C_i} - 1\Bigr),
\end{equation}
which counts turns as well as lengthening, but inherits the noise of inflection detection.

\textbf{B. Turning angles (tangents only).}
The \emph{sum of angles metric} (SOAM) adds the 3D angle between consecutive tangents, in-plane and
out-of-plane, and divides by $L$ \cite{bullitt2003tortuosity}; for fine sampling it equals
$(1/L)\int\lvert\kappa\rvert\,\mathrm ds$. It was the best of Bullitt's metrics at detecting tight,
low-amplitude coils, and the \emph{inflection count metric} caught two other abnormal types.
Keeping the turning angle as a \emph{profile} rather than a sum is what a recent angiography study
did: in 22,334 patients, a transformer reading the full local profile predicted CAD better than its
global mean (AUROC 0.67 against 0.60) \cite{telloayala2026tortuosity}. The clinical scores ---
bends above $45^\circ$, or consecutive $180^\circ$ turns \cite{groves2009tortuosity} --- are
categorical versions of the same idea.

\textbf{C. Curvature integrals (second derivative and above).}
Kashyap et al.\ compared the continuous measures directly against the percentage of vessel area
below 0.4~Pa time-averaged wall shear stress \cite{kashyap2022tortuosity}:
\begin{align}
  \bar\kappa_{a} &= \frac{1}{L}\int_0^L \lvert\kappa(s)\rvert\,\mathrm ds
     && R^2 = 0.112,\ p < 0.001, \\
  \kappa_{\mathrm{rms}} &= \Bigl(\frac{1}{L}\int_0^L \kappa(s)^2\,\mathrm ds\Bigr)^{1/2}
     && R^2 = 0.093,\ p = 0.002, \\
  \bar\kappa'_{2} &= \frac{1}{L}\int_0^L \Bigl(\frac{\mathrm d\kappa}{\mathrm ds}\Bigr)^{2}\mathrm ds
     && R^2 = 0.101,\ p = 0.001,
\end{align}
against $R^2 \approx 0$ ($p = 0.865$) for $L/D$. The mean absolute curvature weights bends
linearly. The squared-derivative measure responds to how abruptly bends start and stop --- plausibly
where flow separates --- but is a third derivative and the noisiest of all. Torsion and
non-planarity (\Cref{eq:week3-frenet}) are the only measures here that recognise a helix as a
helix.

\textbf{The curvature-squared integral.} The measure proposed this week,
\begin{equation}
  T_\kappa = \frac{1}{L}\int_0^L \kappa(s)^2\,\mathrm ds,
  \label{eq:week3-tkappa}
\end{equation}
is not new. It is exactly how Hart et al.\ defined retinal vessel tortuosity (``total squared
curvature'' normalised by length), with 91--95\% classification accuracy
\cite{hart1999tortuosity}, and it is the square of Kashyap's RMS curvature. It is the bending
energy of an elastic rod. Because it weights curvature quadratically, one sharp kink scores higher
than a gentle arc with the same total turning, which is the one thing it adds over
$\bar\kappa_a$. Its unit is mm$^{-2}$, so it is not scale-free; $L\int\kappa^2\,\mathrm ds$ is
dimensionless and equals $\theta^2$ for a circular arc turning through $\theta$.

The ``Jacobian'' framing --- treat the tangent as a flow $\mathbf f$ and take the curvature of its
flow lines, $\kappa = \lvert \mathbf f \times (J\mathbf f) \rvert / \lvert \mathbf f\rvert^3$ ---
reduces on a curve to $(J\mathbf f)\mathbf f = \mathrm d\mathbf T/\mathrm ds$, i.e.\ to
\Cref{eq:week3-frenet}. It needs no new machinery; only the squaring and averaging are new.

\textbf{Its weakness is noise, and squaring makes it worse.} With independent position noise
$\varepsilon$ per coordinate and spacing $h$, the finite-difference second derivative carries noise
of variance $6\varepsilon^2/h^4$ per coordinate. Counting the two directions normal to the tangent,
$\kappa^2$ picks up a \emph{positive} bias of about $12\varepsilon^2/h^4$ that does not average
out. On the tracer's 0.5~mm grid, with $\varepsilon \approx h/\sqrt{12} \approx 0.14$~mm, that is of
order 1--4~mm$^{-2}$, against real coronary values of about 0.003--0.09~mm$^{-2}$
($\kappa \approx 0.05$--0.3~mm$^{-1}$). Raw tracer output would measure the grid: the tracer can
turn only by multiples of about $35^\circ$ or more between 0.5--0.87~mm steps. Gaussian smoothing at
scale $\sigma$ shrinks the bias roughly as $\sigma^{-5}$, while the bias in $\bar\kappa_a$ shrinks
as $\sigma^{-2.5}$. So $T_\kappa$ gains more from smoothing, and also depends more on the choice of
$\sigma$, which must be stated with every value.

\textbf{D. How $\kappa$ is estimated on discrete points.} This choice matters more than which
integral is taken.
\begin{itemize}
  \item \emph{Gaussian smoothing, then finite differences}: the current script ($\sigma = 1$~mm on
    top of the dataset's 0.5~mm).
  \item \emph{Penalised smoothing spline}: its penalty is itself $\int\lvert\mathbf r''\rvert^2$,
    the bending energy, and generalised cross-validation lets the data choose the smoothing scale.
    AneuRisk used free-knot regression splines on carotid centerlines before differentiating
    \cite{sangalli2009aneurisk}.
  \item \emph{Menger curvature}: the curvature of the circle through $\mathbf p_{i-k}$,
    $\mathbf p_i$, $\mathbf p_{i+k}$, $\kappa = 4A/(abc)$ with $A$ the triangle's area and $a,b,c$
    its sides. The scale is explicit: $\kappa$ is measured over $2kh$~mm.
  \item \emph{Discrete bending energy}: following the discrete elastic rods of Bergou et al.\
    \cite{bergou2008rods}, the curvature at vertex $i$ is $2\tan(\varphi_i/2)/\bar\ell_i$,
    with $\varphi_i$ the turning angle and $\bar\ell_i$ the curve length assigned to the vertex
    (half of each adjacent edge), so
    $\int\kappa^2\,\mathrm ds \approx \sum_i \bigl(2\tan(\varphi_i/2)\bigr)^2/\bar\ell_i$. It
    builds $T_\kappa$ from turning angles alone, and is a cross-check on the gradient estimate.
\end{itemize}

\textbf{The field version of the Jacobian idea.} The Tubular Curvature Filter needs no centerline:
it takes the vessel direction field $\mathbf v(\mathbf x)$ from the eigenvectors of the image
Hessian and measures curvature at every voxel from how that field changes along itself
\cite{sunger2023tcf}, in effect $\kappa(\mathbf x) = \lvert (\nabla\mathbf v)\,\mathbf v \rvert$.
This is unchanged when $\mathbf v$ is replaced by $-\mathbf v$, so the eigenvector's arbitrary sign
does not matter, but $\mathbf v$ must be made continuous between neighbouring voxels before it is
differentiated. It avoids the tracer grid entirely and separates inner- from outer-wall curvature;
the smoothing decision moves into the Hessian scale. It is worth considering for CGPS, where only
predicted masks exist.

\begin{figure}[htbp]
  \centering
  \includegraphics[width=\linewidth]{content/background/figures/tortuosity_ambiguity.pdf}
  \caption{Six synthetic centerlines with the same chord and arc length, so the same tortuosity
    index $L/D = 1.25$, coloured by local curvature (\texttt{scripts/make\_tortuosity\_ambiguity.py}).}
  \label{fig:week3-tortuosity-ambiguity}
\end{figure}

\textbf{No single number separates all six curves} (\Cref{tab:week3-comparison}). The proximal and
distal bends are mirror images, so every integral --- $L/D$, SOAM, $\bar\kappa_a$, $T_\kappa$,
$\bar\kappa'_2$ --- gives them identical values. Only a profile along the vessel, or a
position-weighted measure, tells them apart. This agrees with the review's point that the sign of a
tortuosity finding depends on the level of analysis \cite{shen2026geometry}, and with the
profile beating its mean in the angiography study \cite{telloayala2026tortuosity}.

\begin{table}[htbp]
  \centering
  \small
  \begin{tabular}{@{}llllll@{}}
    \toprule
    & & Noise on & \multicolumn{3}{c}{Separates, in \Cref{fig:week3-tortuosity-ambiguity}} \\
    \cmidrule(l){4-6}
    Measure & Derivative & raw tracer & C / S / many & proximal / distal & helix \\
    \midrule
    $L/D$ & 0 & very low & no & no & no \\
    tortuosity density & 0 (+ infl.) & medium & yes & no & no \\
    SOAM, $\bar\kappa_a$ & 1--2 & high & yes & no & no \\
    $T_\kappa$, $\kappa_{\mathrm{rms}}$ & 2, squared & very high & yes$^{*}$ & no & no \\
    $\bar\kappa'_2$ & 3 & extreme & yes & no & no \\
    torsion, non-planarity & 3 & extreme & -- & no & yes \\
    $\kappa(s)$ profile & 2 & high & yes & yes & with $\tau(s)$ \\
    \bottomrule
  \end{tabular}
  \caption{How each measure behaves on unsmoothed tracer output, and which of the curves of
    \Cref{fig:week3-tortuosity-ambiguity} it separates. $^{*}$Weighting sharp bends more than
    gentle ones.}
  \label{tab:week3-comparison}
\end{table}

\textbf{Recommendation for the hand-crafted route.} Compute one smoothed $\kappa(s)$ profile per
segment at a stated scale. Report from it $\bar\kappa_a$ as the primary measure (the best coronary
evidence), $T_\kappa$ or $\kappa_{\mathrm{rms}}$ alongside as the sharp-bend-weighted companion,
and torsion or non-planarity for the out-of-plane part. Keep the profile itself for localisation.
Before trusting $T_\kappa$ on tracer output, measure its noise floor on a straight line quantised to
the 0.5~mm grid.

% ---------------------------------------------------------------------------------------------
\subsection*{5. Letting a network find the structure}

The alternative to choosing the numbers is to let a network choose them. Compress each vessel, or
each stretch of one, into a short vector --- an \emph{embedding} --- without telling the network
what tortuosity is. Then ask whether tortuosity shows up in the embedding, and what else does.

\textbf{The unit decides what can be found} (\Cref{tab:week3-units}). With 800 usable cases, deep
models at the level of the whole tree will largely rediscover PCA. Windows along the paths give
enough samples to train properly.

\begin{table}[htbp]
  \centering
  \small
  \begin{tabular}{@{}lll@{}}
    \toprule
    Unit & Samples & Structure it can reveal \\
    \midrule
    window of 10--30~mm & tens of thousands & bend types, and where each sits \\
    one vessel (LAD, LCx, RCA) & about 800 each & vessel-level shape modes \\
    whole tree & 800 & branching and geometry together \\
    \bottomrule
  \end{tabular}
  \caption{Units of analysis for representation learning, with the number of samples each gives
    in the 800 usable ImageCAS-X cases.}
  \label{tab:week3-units}
\end{table}

\textbf{How the centerline is fed in} matters more than the architecture.
\begin{itemize}
  \item \emph{3D points}, resampled to a fixed number and aligned. Simple, but the network spends
    capacity on pose, and a point-distance reconstruction loss favours the large-scale course and
    blurs small coils, which are exactly the high-frequency part of tortuosity.
  \item \emph{Intrinsic signals} $\kappa(s)$, $\tau(s)$, $r(s)$. Invariant to pose by
    construction, with tortuosity directly in the input; the AneuRisk study used radius and
    curvature profiles in exactly this way \cite{sangalli2009aneurisk}. It inherits the derivative
    noise of \S4.
  \item \emph{Square-root velocity function}, $\mathbf q = \mathbf r'/\sqrt{\lvert\mathbf r'\rvert}$.
    First derivatives only, and its elastic metric separates bending from stretching. It has been
    extended to tree-shaped objects, with branch sliding and branches appearing and disappearing
    \cite{wang2024elastictrees}.
  \item \emph{Lumen surface} as a point cloud or mesh, which carries the radius, so a bend relative
    to vessel width comes for free \cite{beetz2025pointvae, griffo2026wss}.
  \item \emph{Tree graph}, bifurcations as nodes and segments as edges with geometric features ---
    the natural input for whole-tree models.
\end{itemize}

\textbf{Families of methods.}
\begin{description}
  \item[(a) PCA and functional PCA --- the baseline.] Resample, align, find the directions of
    greatest variation, and plot each mode at $\pm 2$ standard deviations to see what it encodes.
    This has been done for the left main bifurcation \cite{medranogracia2017bifurcation} and on
    carotid curvature profiles \cite{sangalli2009aneurisk}. If a network cannot beat it, the
    network adds nothing.
  \item[(b) Elastic shape analysis.] Not a neural network, but the principled middle ground: a
    proper distance between shapes, a geodesic mean and PCA in shape space, with the alignment
    along the vessel handled correctly \cite{wang2024elastictrees}.
  \item[(c) Autoencoders and VAEs on one vessel.] A 1D CNN or transformer compresses the point or
    $\kappa/\tau/r$ sequence to 4--16 numbers and a decoder rebuilds it. A $\beta$-VAE adds
    pressure for each latent dimension to capture one factor \cite{higgins2017betavae}, read off by
    \emph{latent traversal}: vary one dimension, decode, watch what changes. In cardiac anatomy,
    individual components of such models mapped onto recognisable 3D shape features, and their
    encodings predicted myocardial infarction better than clinical benchmarks
    \cite{beetz2022meshvae, beetz2025pointvae}.
  \item[(d) Self-supervised learning on windows --- the best fit here.] \emph{Contrastive:} two
    altered copies of a window are pulled together in embedding space and pushed away from other
    windows. The alterations define what counts as noise --- random rotation, resampling jitter,
    quantisation to the 0.5~mm grid, a changed smoothing $\sigma$ --- so the embedding is trained
    to ignore exactly the artefacts that make hand-crafted curvature unstable, which also serves
    the reproducibility objective. \emph{Masked:} hide a stretch of centerline and predict it from
    its surroundings; to do so the network must learn how vessels bend. Clustering the window
    embeddings gives a bend vocabulary; colouring each tree by cluster shows where each bend type
    occurs.
  \item[(e) Whole-tree models.] Graph neural networks already capture coronary topology and
    position well, though they have been trained to label segments
    \cite{hampe2024gnn, zhao2023eagmn}; replacing the labelling head with a self-supervised
    objective, and using the window embeddings of (d) as edge features, gives a tree embedding.
    A recursive VAE encodes the binary vessel tree bottom-up and reproduced real tortuosity with a
    similarity of 0.96 \cite{feldman2023vesselvae}. Part-based generation separates topology from
    segment geometry explicitly \cite{chen2025partvessel}.
  \item[(f) Topological data analysis.] Persistence diagrams summarise branching at multiple
    scales; on brain artery trees they correlated with age and sex more strongly than earlier
    analyses \cite{bendich2016persistent}. A non-deep comparison at the tree level.
\end{description}

\textbf{How to tell whether the network found tortuosity.}
\begin{enumerate}
  \item \emph{Linear probes}: regress each hand-crafted measure ($L/D$, $\bar\kappa_a$, $T_\kappa$,
    torsion, inflections) on the embedding and report $R^2$.
  \item \emph{The residual}: regress the hand-crafted panel out of the embedding and inspect what
    remains by latent traversal or cluster maps --- structure nobody thought to measure.
  \item \emph{The six curves} of \Cref{fig:week3-tortuosity-ambiguity}: a good representation
    separates all of them, the proximal/distal pair only if position along the vessel is kept.
  \item \emph{Invariance}: embed the same vessel at different smoothing and sampling and compare
    that spread with the spread between vessels (an intraclass correlation).
  \item \emph{Beat PCA} on reconstruction and on the probes.
  \item \emph{Downstream}, on CGPS: does the embedding add to the hand-crafted measures for plaque
    progression?
\end{enumerate}
t-SNE or UMAP are for looking, not measuring.

\textbf{Pitfalls specific to this setting.}
\begin{itemize}
  \item \emph{The network learns the pipeline, not the anatomy.} On CGPS the centerlines come from
    predicted masks, and minimal-path centerline methods ``struggle with tortuosity and often tend
    to make shortcuts'' \cite{korona2025tracing}, which would bias learned tortuosity low on the
    most tortuous vessels. Scanner, image-quality grade (1--4 in the usable ImageCAS-X cases) and
    segmentation breaks can all be encoded; check embeddings against quality grade first.
  \item \emph{Pose, length and sampling density} dominate raw coordinates unless removed by the
    input representation or the augmentations.
  \item \emph{Point-distance losses blur small bends}, as noted above.
  \item \emph{Interpretability.} Embedding dimensions have no units; a clinical reader will still
    want mm$^{-1}$, and the probes are what connect the two.
\end{itemize}

\textbf{Suggested path.}
(1) PCA or functional PCA per segment for LAD, LCx and RCA.
(2) A self-supervised encoder on windows of $\kappa/\tau/r$ or point sequences, with
measurement-nuisance augmentations; cluster into a bend vocabulary and map it onto the trees.
(3) Probe the embeddings against the hand-crafted panel and inspect the residual.
(4) Only then a tree-level GNN or recursive VAE, with the window embeddings as edge features.

% ---------------------------------------------------------------------------------------------
\subsection*{6. Open questions}

\begin{itemize}
  \item Whole-vessel tortuosity needs each label's pieces chained from its origin to its tip; the
    current per-segment rows describe only the longest piece.
  \item Which smoothing scale $\sigma$ is defensible, and is it the same for ImageCAS-X centerlines
    (already smoothed at 0.5~mm) and for tracer output on CGPS? A quantised-synthetic-curve test
    would settle the noise floor of each measure.
  \item For Bj{\o}rn: is a local bend vocabulary the kind of ``structure'' the thesis wants, or is a
    per-vessel or per-tree embedding more useful for the CGPS risk model?
\end{itemize}
